\medskip

On veut déterminer la forme exponentielle du nombre complexe 
$z = - 6\sqrt 3 + 6\text{i}.$

\begin{list}{\textbullet}{}
\item Le module de $z$ est $|z| = \ds\sqrt{\left(-6\sqrt{3}\right )^2 + 6^2 } = \ds\sqrt{108+36} = \ds\sqrt{144}=12$.

\item Un argument $\theta$ de $z$ vérifie $\cos\left (\theta\right )=\dfrac{-6\sqrt{3}}{|z|} = \dfrac{-6\sqrt{3}}{12} = -\dfrac{\sqrt{3}}{2}$ et $\sin\left (\theta\right )=\dfrac{6}{|z|} = \dfrac{6}{12} = \dfrac{1}{2}$.

$\theta = \dfrac{5\pi}{6}$ vérifie ces égalités.
\end{list}

La forme exponentielle du nombre complexe 
$z = - 6\sqrt 3 + 6\text{i}$
est donc $z = 12 \e^{\frac{5\pi}{6}\i}$.

\bigskip
